\scanpage{37}
pour $n$ variable. Ceci dit, on a un hom évident
\diagram{\[
(41)\qquad\begin{tikzcd}[ampersand replacement=\&,column sep=huge]
\mathcal E=(\unsure{\mathfrak S_3|\mathbb G_3}\times\Gamma)
\mathbin{\underset{\scriptscriptstyle1/2}{.}}\pi_\infty\arrow[r]\&
\widetilde{\mathcal E}=(\unsure{\mathfrak S_3|\mathbb G_3}\times A)
\mathbin{\underset{\scriptscriptstyle1/2}{.}}\pi_\infty
\end{tikzcd}
\]}
déduit de (40), et les opérations de $\mathcal E$, sur les ens. précédents
se font \uline{à travers} $\widetilde{\mathcal E}\ldots$

(1.9) Les développements précédents n'utilisaient pas que la clôture alg.
\inserted{envisagée} $\overline{\mathbb Q}$ de $\mathbb Q$ était plongée dans
$\mathbb C$ -- ils marcheraient $\unsure{=|\text{en}}$ en remplaçant
$\mathbb Q$ par n'importe quel corps premier, à condition de se borner à des
$n$ premiers à la car. ($\overline C_\infty$ est alors le revêtement universel
abélien \uline{modéré} de $(\overline C_1,\unsure{\overline I|I})\ldots$).
Il faut l'utiliser, pour faire le lien avec les cartes 3-pondérées finies
abéliennes, et le groupe cartographique. On va d'abord \unsure{pointer|prendre}
$\overline C_\infty$ au dessus de
$\alpha=-\overline j\in\omega=\{\alpha,\overline\alpha\}$, pour préciser
l'identification au rev. universel profini \deleted{\unsure{pointé}} \uline{en $\alpha$}.
% archive foliation: 36
